% [原文 p.105]
\begin{proof}[证明（续）]
同样地，对 $n>n_0$，$L$ 与 $D(nL,H):=na(K)L+a(K)H+r_L(H)K$ 都是非丰富的 nef 除子。因此 $0\neq F_{D(nL,H)}\subset F_L$，且 $F_{D(nL,H)}\subset\overline{NE}_{K<0}\cup\{0\}$ 成立。

\par\medskip\noindent\textbf{步骤 2}\quad 下面证明：若 $\dim F_L>1$ 且 $F_L\not\subset\overline{NE}_{K\geq 0}$，则可适当选取丰富除子 $H$，使
\[
\dim F_{D(nL,H)}<\dim F_L.
\]
为此，选取丰富除子 $H_i$，使它们在 $F_L$ 上的限制构成 $F_L^*$ 的基底。若 $\dim F_L>1$，则线性函数
\[
\left.\left(nL+H_i+\frac{r_L(H_i)}{a(K)}K\right)\right|_{F_L}=\left.\left(H_i+\frac{r_L(H_i)}{a(K)}K\right)\right|_{F_L}
\]
不会在 $F_L$ 上恒等于 $0$。因此，对适当的 $i$ 有 $\dim F_{D(nL,H_i)}<\dim F_L$，断言得证。

把这一讨论依次应用于更小的锥面，即可知：给定满足 $F_L\not\subset\overline{NE}_{K\geq 0}$ 的任意 nef 除子 $L$，可以适当选取 nef 除子 $L'$，使 $F_L\supset F_{L'}$、$\dim F_{L'}=1$、$F_{L'}\subset\overline{NE}_{K<0}\cup\{0\}$ 同时成立。

\par\medskip\noindent\textbf{步骤 3}\quad 证明下述等式：
\[
\overline{NE}=\overline{\left(\overline{NE}_{K\geq 0}+\sum_{\dim F_L=1}F_L\right)}.
\]
为此假定右端较小。此时存在除子类 $M\in N_{\mathbb{Z}}^*$，使超平面 $M=0$ 与右端不相交，而与左端相交。于是可以作与定理 1.24 相同的讨论。应用有理性定理与步骤 2，即得矛盾。
% [Tutor Note: 步骤 3 的等式表明曲线锥由 $(K_X+\Delta)$ 非负部分与一维极锥面拼成；剩下的困难只在于证明这些一维 $F_L$ 由曲线生成，因此步骤 4 转向极射线在半空间 $K<0$ 中的聚积性。]

\par\medskip\noindent\textbf{步骤 4}\quad 其次证明一维的 $F_L$ 不在半空间 $K<0$ 中聚积。为此，选取丰富除子 $H_1,\dots,H_d\in N_{\mathbb{Z}}^*$，使 $K$ 与 $H_1,\dots,H_d$ 构成 $N_{\mathbb{R}}^*$ 的基底。

% [原文 p.106]
对各一维的 $F_L$ 与 $i$，取 $n_i$ 使 $F_{D(n_iL,H_i)}=F_L$；于是对生成 $F_L$ 的 $\xi$ 以及所有的 $i$，
\[
\frac{(\xi\cdot H_i)}{(\xi\cdot K)}=\frac{\text{整数}}{a(K)}
\]
成立。若 $F_L$ 在半空间 $K<0$ 的某处聚积，则在射影化 $\mathbb{P}(N_{\mathbb{R}})$ 中，与它们对应的点就在由 $K<0$ 给出的仿射子集 $U\subset\mathbb{P}(N_{\mathbb{R}})$ 的某处聚积。但由上述等式可知这不可能，因为
\[
\xi\in U\mapsto\frac{(\xi\cdot H_1)}{(\xi\cdot K)},\dots,\frac{(\xi\cdot H_d)}{(\xi\cdot K)}
\]
给出了仿射坐标系。

由此可知，满足 $F_L\cdot(K+\epsilon H)<0$ 的 $F_L$ 至多只有有限个。因此
\[
\overline{NE}_{K+\epsilon H\geq 0}+\sum_{F_L\cdot(K+\epsilon H)<0}F_L
\]
是闭锥。步骤 3 已证明其闭包等于 $\overline{NE}$。于是
\[
\overline{NE}=\overline{NE}_{K+\epsilon H\geq 0}+\sum_{F_L\cdot(K+\epsilon H)<0}F_L.
\]
至此，除去「各 $F_L$ 由定理 3.7 中所说的曲线生成」这一断言，定理 3.7(2) 已获证明。

\begin{remark}{3.16}
行文之便，这里暂且离开主线，先证明补充 3.6。锥定理的证明从步骤 5 继续。

为避免记号混乱，设 $M$ 是补充 3.6 中 nef 且大的除子。

适当取 $m_i\in\mathbb{Q}$，可把 $M$ 表为 $M\equiv m_0K_X+\sum m_iH_i$。步骤 4 已看到：当 $\xi$ 取遍所有 $(K_X+\Delta)$-负的极射线时，分数 $-(\xi\cdot M)/(\xi\cdot(K_X+\Delta))$ 的分母有界。因此该集合的下确界不外是它的最小值，由某条极射线 $\xi_M$ 达到。于是得到 $r(M)=-(\xi_M\cdot M)/(\xi_M\cdot(K_X+\Delta))$ 与 $\xi_M\cdot(M+r(M)\cdot(K_X+\Delta))=0$。\hfill$\blacksquare$
\end{remark}

% [原文 p.107]
\par\medskip\noindent\textbf{步骤 5}\quad 定理 3.15 即定理 3.7(1) 是定理 3.7(2) 与步骤 4 的形式推论，这一点可由定理 1.24 证明最后部分的讨论说明。

\par\medskip\noindent\textbf{步骤 6}\quad 这里证明：若 $F\subset\overline{NE}(X)$ 是 $(K_X+\Delta)$-负的极锥面，则存在 nef 的 Cartier 除子 $D$，使 $F_D=F$。

设 $\langle F\rangle$ 是 $F$ 生成的子向量空间，$V\subset N_1(X)^*$ 是在 $\langle F\rangle$ 上取值为 $0$ 的线性函数全体之集。$F$ 由极射线生成，而各极射线都在 $\mathbb{Q}$ 上有定义，故 $V$ 也在 $\mathbb{Q}$ 上有定义。选取 $\epsilon>0$，使 $K_X+\Delta+\epsilon H$ 在 $F$ 上为负。

由于 $F$ 是极端的，$\langle F\rangle\cap\overline{NE}(X)=F$。因此
\[
W_F:=\overline{NE}(X)_{K_X+\Delta+\epsilon H\geq 0}+\sum_{\dim F_L=1,\,F_L\not\subset F}F_L
\]
是闭锥，与 $\langle F\rangle$ 只交于原点，并具有 $\overline{NE}(X)=W_F+F$ 的性质。因此，存在超平面 $(g=0)$，它包含 $\langle F\rangle$，但与 $W_F$ 只交于原点。

$F$ 的支撑函数（supporting function）不是别的，正是 $V$ 中的元素，且它在 $W_F\setminus\{0\}$ 上严格为正（$>0$）。这些元素全体构成开集，又包含 $g$，故非空。由此可知 $F$ 具有在 $\mathbb{Q}$ 上有定义的支撑函数。把它适当乘上若干倍，即得对应于某个 Cartier 除子（记作 $D$）的支撑函数。

\par\medskip\noindent\textbf{步骤 7}\quad 走到这一步，定理 3.7(3) 便是无基点定理的简单推论。

由假设，$-(K_X+\Delta)$ 在 $F$ 上严格为正。这意味着对 $m\gg 1$，$mD-(K_X+\Delta)$ 在 $\overline{NE}(X)\setminus\{0\}$ 上严格为正。因此由定理 1.18，$m\gg 1$ 时 $mD-(K_X+\Delta)$ 是丰富的。由此从定理 3.3 可知，$m\gg 1$ 时 $|mD|$ 无基点。对适当的 $m\gg 1$，把由 $|mD|$ 决定的态射的 Stein 分解 [Har77, III.11.5] 记作 $g_F$。它不外是附属于 $F$ 的收缩态射。

\par\medskip\noindent\textbf{步骤 8}\quad 若 $F\neq\{0\}$，则由定理 1.18 可知相应的收缩态射不是同构。因此至少有一条曲线 $C\subset X$ 被 $g_F$ 收缩。这意味着每个 $(K_X+\Delta)$-负的极锥面都含有曲线的类。

% [原文 p.108]
若 $X$ 非奇异，则用定理 1.13 可以证明 $F$ 内存在有理曲线。（这留作习题。）一般情形在 [Kaw91] 中用与定理 1.10 类似的方法证明。

\par\medskip\noindent\textbf{步骤 9}\quad 设 $F$ 是 $(K_X+\Delta)$-负的极锥面。设 $D$ 是 $F$ 的任意支撑 Cartier 除子，$g_F:X\to Z$ 是 $F$ 的收缩态射，$W_F\subset\overline{NE}(X)$ 如步骤 6 所述。如步骤 7 所示，$g_F$ 由任意 $m\gg 1$ 的 $|mD|$ 给出。因此 $mD$ 与 $(m+1)D$ 都是 $Z$ 上某个 Cartier 除子的拉回。取差即得满足 $g_F^*D_Z=D$ 的 Cartier 除子 $D_Z$。

设 $L$ 是满足 $L|_F\equiv 0$ 的 $X$ 上任意 Cartier 除子。$D$ 在 $W_F\setminus\{0\}$ 上为正，故 $m\gg 1$ 时 $L+mD$ 也在 $W_F\setminus\{0\}$ 上为正。于是 $L+mD$ 是 $F$ 的支撑 Cartier 除子；又如上所示，存在 Cartier 除子 $M_Z$ 使 $L+mD=g_F^*M_Z$。令 $L_Z:=M_Z-mD_Z$，即得定理 3.7(4)。
\end{proof}
% [Tutor Note: 补充 3.6 的证明只用到步骤 4 的定量结论：负极射线上分数 $-(\xi\cdot M)/(\xi\cdot(K_X+\Delta))$ 的分母有界，故其下确界可由某条极射线达到。这正是「临界斜率处必出现负极射线」的精确表述。]

接下来的两个推论与收缩态射的除子论有关。

\begin{corollaryp}{3.17}
$(X,\Delta)$ 是射影 klt 对，$R\subset\overline{NE}(X)$ 是 $(K_X+\Delta)$-负的极射线，$g_R:X\to Z$ 是其收缩态射。设 $C\subset X$ 是生成 $R$ 的曲线。此时有下述正合列：
\[
0\to\Pic Z\xrightarrow{L\mapsto g_R^*L}\Pic X\xrightarrow{M\mapsto(M\cdot C)}\mathbb{Z}.
\]
特别地，$\rho(Z)=\rho(X)-1$。
\end{corollaryp}

\begin{proof}
设 $L$ 是 $Z$ 上的线丛。此时 $(g_R)_*(g_R^*L)=L\otimes(g_R)_*\mathcal{O}_X=L$ 成立，故 $L\mapsto g_R^*L$ 是单射。若 $M$ 是 $X$ 上的线丛且满足 $(M\cdot C)=0$，则由定理 3.7(4) 存在 $L$ 使 $M=g_R^*L$。
% [Tutor Note: 这一行原书印作 $(N\cdot C)=0$，但本段讨论的线丛是 $M$，$N$ 未在上下文出现；英文底本亦作 $(M\cdot C)=0$。译文按数学含义写作 $(M\cdot C)=0$，此处留痕备查。]
\end{proof}
% [Tutor Note: 正合列中段给出「线丛沿纤维为平凡则可下降」，末段用相交数判定。$\rho(Z)=\rho(X)-1$ 说明收缩一条极射线恰使 Picard 数减少 1，与相对极小模型的每次收缩一致。]

\begin{corollaryp}{3.18}
$(X,\Delta)$ 是射影 klt 对，$R\subset\overline{NE}(X)$ 是 $(K_X+\Delta)$-负的极射线，$g_R:X\to Z$ 是其收缩态射。设 $X$ 是 $\mathbb{Q}$-分解的，且 $g_F$ 是除子收缩或 Fano 收缩（2.14 步骤 3(1)）。此时 $Z$ 也是 $\mathbb{Q}$-分解的。
\end{corollaryp}

% [原文 p.109]
\begin{proof}
首先设 $g_R$ 是除子收缩，$E\subset X$ 是例外除子。特别地，$(E\cdot R)\neq 0$。设 $B$ 是 $Z$ 上的 Weil 除子。此时存在有理数 $s$，使 $((g_R)_*^{-1}B+sE)\cdot R=0$。适当取 $m>0$，则 $m((g_R)_*^{-1}B+sE)$ 是 Cartier 除子，且由定理 3.7(4) 可知它是 $Z$ 上某个 Cartier 除子 $M_Z$ 的拉回。由此 $mB\sim M_Z$，故 $B$ 是 $\mathbb{Q}$-Cartier。

其次设 $g_R$ 是 Fano 收缩。设 $B$ 是 $Z$ 上的不可约 Weil 除子，$Z^0\subset Z$ 是非奇异点全体构成的子集。设 $D\subset X$ 是 $g_F^{-1}(B|_{Z^0})$ 的闭包。$D$ 不与 $g_F$ 的一般纤维相交，故 $(D\cdot R)=0$。适当取 $m>0$，则 $mD$ 是 Cartier 除子，且由定理 3.7(4) 可知它是 $Z$ 上某个 Cartier 除子 $M_Z$ 的拉回。于是 $mB\sim M_Z$，故 $B$ 是 $\mathbb{Q}$-Cartier。
\end{proof}
% [Tutor Note: 原文在此处把收缩态射写作 $g_F$（$F$ 指一般纤维），与前面定义的 $g_R$ 是同一个态射；后文 §3.4 也用 $g$ 与 $g_F$ 交替，读者按语境辨认即可。]

\section{有理性定理}
\takeaway{有理性定理断言临界值 $r(H)$ 是有理数且分母有界；证明把一般情形归约为 $H$ 无基点的情形，再用关于格点分布的数论引理与 Noether 式归纳法导出矛盾。}

\begin{proof}
本节证明取自 [Kaw84a]，并按 [Kol84] 作了简化与补充。（参看 [KMM87, 4.1]。）

主要想法如下。为简单起见，假定 $X$ 非奇异、$H$ 丰富且 $\Delta=0$。我们想用无基点手法证明某个形如 $|pH+qK_X|$ 的线性系无基点。仔细检查证明可知，本质之处在于：消没定理 $H^1(X,\mathcal{O}_X(pH+qK_X))=0$ 成立，恰在 $pH+(q-1)K_X$ 为 nef 且大之时。若选取 $p,q$ 使
\begin{equation}
\frac{q-1}{p}<r<\frac{q}{p}
\tag{3.1}
\end{equation}
成立，则上述消没定理成立，但 $|pH+qK_X|$ 连 nef 都不是，更谈不上无基点。当然还需要其他消没定理，所以能证明的只是：若满足不等式 (3.1) 的对 $(p,q)$ 足够多，则可导出矛盾。步骤 3 将说明，这种大量对的存在与 $r$ 的数论性质如何相互关联。

\par\medskip\noindent\textbf{步骤 1}\quad 本步骤把定理 3.5 的一般情形归结为 $H$ 无基点这一特殊情形。假定 $a(K_X+\Delta)$ 是 Cartier 的。由定理 3.3，对 $m\gg c\gg d>0$，$H':=m(cH+da(K_X+\Delta))$ 无基点。解线性方程
% [原文 p.110]
\[
H+r(H)(K_X+\Delta)=\lambda\left(H'+r(H')(K_X+\Delta)\right)
\]
得
\[
r(H)=\frac{r(H')+mda}{mc}.
\]
因此，$r(H)$ 是有理数，恰在 $r(H')$ 是有理数之时。进而，设 $r(H')$ 是有理数，其分母为 $v$。此时 $r(H)$ 的分母整除 $mcv$。由于 $m,c$ 可取任意大，这意味着 $r(H)$ 的分母整除 $v$。

\par\medskip\noindent\textbf{步骤 2}\quad 设 $Y$ 是非奇异射影簇，$\{D_i\}$ 是 Cartier 除子的有限集合，$A$ 是正规交叉 $\mathbb{Q}$-除子，且 $\lceil A\rceil$ 有效。考虑 Hilbert 多项式
\[
P(u_1,\dots,u_k)=\chi\left(\sum u_iD_i+\lceil A\rceil\right).
\]
假定对 $u_i$ 的某个值，$\sum u_iD_i$ 是 nef 的，而且 $\sum u_iD_i+A-K_Y$ 是丰富的。

此时，对所有整数 $m\gg 0$，$\sum mu_iD_i+A-K_Y$ 仍是丰富的，故由消没定理，$H^i(\sum mu_iD_i+\lceil A\rceil)=0$ 对 $i>0$ 成立；又由非消没定理，$\mathcal{O}_Y(\sum mu_iD_i+\lceil A\rceil)$ 有非零截面。特别地，$\chi(\sum mu_iD_i+\lceil A\rceil)\neq 0$。

因此在这种情况下，次数至多为 $\dim Y$ 的多项式 $P(u_1,\dots,u_k)$ 不恒等于 $0$。

\par\medskip\noindent\textbf{步骤 3}\quad 需要关于多项式解集的下述结果。

\begin{lemmap}{3.19}
给定正整数 $a$、正数 $\epsilon$ 与实数 $r$。设 $P(x,y)$ 是次数至多为 $n$ 的非零多项式，并设方程 $0<ay-rx<\epsilon$ 的所有充分大的整数解 $(x,y)$ 都满足 $P(x,y)=0$。此时 $r$ 是有理数，且 $0<\text{分母}\leq a(n+1)/\epsilon$。
\end{lemmap}

% [原文 p.111]
\begin{proof}
首先假定 $r$ 是无理数。于是，在直线 $ay-rx=0$ 的任何一侧，都有无限多个 $(x,y)$-平面的格点（坐标为整数值的点）与该直线的距离小于 $\epsilon/(n+2)$。特别地，存在 $0<ay'-rx'<\epsilon/(n+2)$ 的充分大的整数解 $(x',y')$。因此，由假定，$(x',y'),\dots,((n+1)x',(n+1)y')$ 都是 $P=0$ 的整数解。$P$ 与 $(y'x-x'y)$ 有这 $(n+1)$ 个公共解，故 $(y'x-x'y)$ 整除 $P$。选取更小的 $\epsilon$ 可以把同一讨论重复任意多次，而重复 $(n+1)$ 次即得矛盾。

其次取最简分数表示 $r=u/v$（$v>0$）。任取 $j\in\mathbb{Z}_{>0}$，设 $(x',y')$ 是 $ay-rx=aj/v$ 的解。（整数解的存在容易看出。）此时对所有 $k$ 都有 $a(y'+ku)-r(x'+akv)=aj/v$。因此如上所见，若 $aj/v<\epsilon$，则 $(ay-rx)-(aj/v)$ 整除 $P$。故这样的 $j$ 至多只有 $n$ 个。由此 $a(n+1)/v\geq\epsilon$。
\end{proof}

\par\medskip\noindent\textbf{步骤 4}\quad 设 $\epsilon$ 是正数，$H$ 是 nef 且大的 Cartier 除子。取 $a\in\mathbb{Z}_{>0}$，使 $a(K_X+\Delta)$ 是 Cartier 的。假定 $K_X+\Delta$ 不是 nef 的，并定义
\[
r=\max\{t\in\mathbb{R}:H+t(K_X+\Delta)\ \text{是 nef 的}\}.
\]

对每个 $(p,q)$，$L(p,q)$ 是 $X$ 上线性系 $|pH+qa(K_X+\Delta)|$ 的（带约化结构的）基点集。按定义，$L(p,q)=X$ 与 $|pH+qa(K_X+\Delta)|=\varnothing$ 等价。

\par\medskip\noindent\textbf{主张 3.20}\quad 无论怎样选取满足 $0<aq-rp<\epsilon$ 的充分大的 $(p,q)$，$L(p,q)$ 都是 $X$ 的同一子集。把这个子集记作 $L_0$。

\begin{proof}
让 $X$ 上的除子 $xH+ya(K_X+\Delta)$ 对应于点 $(x,y)$，考虑图 3.1。图 3.1 表示 $L(p',q')\subset L(p,q)$；利用关于子簇的 Noether 式归纳法即可证明断言。
\end{proof}

对上述那样的 $(p,q)$，线性系 $|pH+qa(K_X+\Delta)|$ 不是 nef 的，故在 $X$ 上不可能无基点。设 $I\subset\mathbb{Z}\times\mathbb{Z}$ 是由同时满足 $0<aq-rp<1$ 与 $L(p,q)=L_0$ 的 $(p,q)$ 全体构成的集合。要强调的是，$I$ 包含满足 $0<aq-rp<1$ 的所有充分大的 $(p,q)$。

\par\medskip\noindent\textbf{步骤 5}\quad 设 $(X,\Delta)$ 是 klt 对，$g:Y\to X$ 是对数奇点消解。

% [原文 p.112]
\begin{figure}[htbp]
\centering
\resizebox{0.80\textwidth}{!}{\begin{tikzpicture}[line width=0.5pt]
% 非 nef 区域（$ay-rx>1$）
\fill[gray!25] (-1.30,4.45) -- (1.72,4.45) -- (-1.30,-1.10) -- cycle;
% 三条直线
\draw[dashed] (-1.20,-1.15) -- (1.75,4.50);
\draw (-0.50,-1.15) -- (2.25,4.50);
\draw (-0.32,-1.15) -- (2.42,4.50);
% 坐标轴
\draw[->] (-1.35,0) -- (2.78,0);
\draw[->] (0,-1.05) -- (0,4.45);
\node[anchor=north west] at (2.74,0.04) {$x$};
\node[anchor=south west] at (0.10,4.25) {$y$};
% 直线标号
\node[anchor=east] at (-1.30,-1.00) {$ay-rx=1$};
\node[anchor=west] at (2.44,4.10) {$ay-rx=0$};
% 区域标号（式子在图题中给出，图内只放短标签）
\node[anchor=west] at (-1.26,3.30) {非 nef};
\node[anchor=west] at (1.20,0.60) {丰富};
% 点与箭头
\fill (-0.30,0.25) circle (0.05);
\node[left] at (-0.37,0.25) {$(p,q)$};
\fill (0.05,0.70) circle (0.05);
\node[above] at (0.05,0.76) {$(kp,kq)$};
\draw[->] (0.08,0.71) -- (1.05,1.06);
\fill (1.10,1.08) circle (0.055);
\node[anchor=south west] at (1.14,1.12) {$(p',q')$};
\end{tikzpicture}
}
\figtitle{3.1}{横轴为 $x$，纵轴为 $y$；过原点的斜线是 $ay-rx=0$（斜率 $r/a$），虚线是 $ay-rx=1$。左上方区域中 $xH+ya(K_X+\Delta)$ 不是 nef 的，右下方区域中它丰富；箭头把 $(p,q)$ 移到斜率更小的 $(p',q')$，使对应除子落入丰富区域（原书 p.112）。}
\end{figure}
% [Tutor Note: 页图左上方区域标“非 nef”，右下方区域标“丰富”。原书在图的右侧另有一段说明，意为：当斜率小于丰富的上界再减去 $\varepsilon$ 时，这条箭头延长到一定长度后即成为非常丰富。重绘时该段说明已移入本注，以免图内文字压迫线条。]

设 $D_1=g^*H$、$D_2=g^*(a(K_X+\Delta))$，并按公式 $K_Y\equiv g^*(K_X+\Delta)+A$ 定义 $A$。$\lceil A\rceil$ 是有效的，且是 $g$-例外的。

这里定义 $P(x,y)=\chi(xD_1+yD_2+\lceil A\rceil)$。

$D_1$ 是 nef 且大的，故由 Riemann-Roch 弱定理 1.36 可知 $P$ 不恒等于 $0$。又 $\lceil A\rceil$ 有效且是 $g$-例外的，故
\[
H^0(Y,pD_1+qD_2+\lceil A\rceil)=H^0(X,pH+qa(K_X+\Delta)).
\]

\par\medskip\noindent\textbf{步骤 6}\quad 现在假定有理性定理中「$r$ 是有理数」这一断言不成立。此时若 $0<ay-rx<1$，则 $xD_1+yD_2+A-K_Y$ 数值等价于 nef 且大的 $\mathbb{Q}$-除子 $xH+(ay-1)(K_X+\Delta)$ 的拉回。因此可以应用消没定理，得
\[
H^i(Y,xD_1+yD_2+\lceil A\rceil)=0,\qquad\forall i>0.
\]
由步骤 3，存在任意大的 $(p,q)$，满足 $0<aq-rp<1$ 且 $P(p,q)=h^0(Y,pD_1+qD_2+\lceil A\rceil)\neq 0$。（否则 $P(x,y)$ 将过于频繁地取零，容易推出 $r$ 关于 $X$ 与 $H$ 是有理数。）因此由步骤 4，对所有 $(p,q)\in I$ 都有 $|pH+qa(K_X+\Delta)|\neq\varnothing$。也就是说，$L_0\neq X$。

\par\medskip\noindent\textbf{步骤 7}\quad 对 $(p,q)\in I$，选取对数奇点消解 $f:Y\to X$，使之满足下述条件。

（1）$K_Y\equiv f^*(K_X+\Delta)+\sum a_jF_j$，其中 $a_j>-1$。

（2）$f^*(pH+(qa-1)(K_X+\Delta))-\sum p_jF_j$ 对适当的充分小的 $p_j>0$ 是丰富的（$pH+(qa-1)K_X$ 是 nef 且大的，故由命题 2.61 这是可能的）。

（3）$f^*|pH+qa(K_X+\Delta)|=|L|+\sum r_jF_j$ 对适当的整数 $r_j\geq 0$ 成立。其中 $L$ 无基点，$\sum r_jF_j$ 是固定部分。

% [原文 p.113]
\par\medskip\noindent\textbf{步骤 8}\quad 按步骤 7 的方式选取 $(p,q)\in I$。与前面同样，可以选取有理数 $c>0$ 与 $p_j>0$，使
\[
\sum(-cr_j+a_j-p_j)F_j=A'-F
\]
成立。其中 $\lceil A'\rceil$ 有效，且 $A'$ 中不出现 $F$。比较系数可知，$F$ 被映入线性系 $|pH+qa(K_X+\Delta)|$ 的基点集 $L(p,q)=f(\bigcup_{r_j>0}F_j)$ 的某个成分 $B$ 中。如下定义：
\[
\begin{aligned}
N(p',q')&=f^*(p'H+q'a(K_X+\Delta))+A'-F-K_Y\\
&\equiv f^*\left((p'-(1+c)p)H+(q'-(1+c)q)a(K_X+\Delta)\right)\\
&\quad+f^*\left((1+c)pH+(1+c)qa(K_X+\Delta)\right)\\
&\quad+\sum(-cr_j+a_j-p_j)F_j-K_Y\\
&\equiv cL\ (\text{无基点})\\
&\quad+f^*\left((p'-(1+c)p)H+(q'-(1+c)q)a(K_X+\Delta)\right)\\
&\quad+f^*\left(pH+(qa-1)(K_X+\Delta)\right)-\sum p_jF_j\ (\text{丰富}).
\end{aligned}
\]
注意：若充分大的 $p'$ 与 $q'$ 满足 $aq'-rp'<aq-rp$，则 $(q'-(1+c)q)a<r(p'-(1+c)p)$，且 $(p'-(1+c)p)H+(q'-(1+c)q)a(K_X+\Delta)$ 是 nef 的。由此 $N(p',q')$ 是丰富的。因此由消没定理，下面的映射是满射：
\[
H^0\left(Y,f^*(p'H+q'a(K_X+\Delta))+\lceil A'\rceil\right)\to H^0\left(F,(f^*(p'H+q'a(K_X+\Delta))+\lceil A'\rceil)|_F\right).
\]

\par\medskip\noindent\textbf{步骤 9}\quad 由伴随公式，除子 $f^*(p'H+q'a(K_X+\Delta))+A'-F-K_Y$ 在 $F$ 上的限制由下述除子给出：
\[
(f^*(p'H+q'a(K_X+\Delta))+A')|_F-K_F.
\]

% [原文 p.114]
由步骤 2，$(p',q')$ 的 Hilbert 多项式
\[
\chi\left(F,(f^*(p'H+q'a(K_X+\Delta))+\lceil A'\rceil)|_F\right)
\]
不恒等于 $0$。

当 $0<aq'-rp'<aq-rp$ 时，
\[
(f^*(p'H+q'a(K_X+\Delta))+A')|_F-K_F=N(p',q')|_F
\]
是丰富的，且在该带状区域中下式成立：
\[
\chi\left(F,(f^*(p'H+q'a(K_X+\Delta))+\lceil A'\rceil)|_F\right)=h^0\left(F,(f^*(p'H+q'a(K_X+\Delta))+\lceil A'\rceil)|_F\right).
\]
把步骤 3 应用于 $F$ 的 Hilbert 多项式与 $\epsilon=aq-rp$，可知存在任意大的 $(p',q')$，满足 $0<aq'-rp'<aq-rp$ 与下式：
\[
h^0\left(F,(f^*(p'H+q'a(K_X+\Delta))+\lceil A'\rceil)|_F\right)\neq 0.
\]

\par\medskip\noindent\textbf{步骤 10}\quad 至此已作好导出矛盾的准备。由假定 $L(p,q)=L_0$。对步骤 9 中选取的 $(p',q')$，
\[
H^0\left(Y,f^*(p'H+q'a(K_X+\Delta))+\lceil A'\rceil\right)\to H^0\left(F,(f^*(p'H+q'a(K_X+\Delta))+\lceil A'\rceil)|_F\right)\neq 0
\]
是满射。这意味着 $F$ 不包含在 $|f^*(p'H+q'a(K_X+\Delta))+\lceil A'\rceil|$ 的基点集中。由于 $\lceil A'\rceil$ 是 $f$-例外的且有效，
\[
H^0\left(Y,f^*(p'H+q'a(K_X+\Delta))+\lceil A'\rceil\right)=H^0\left(X,p'H+q'a(K_X+\Delta)\right).
\]
如 3.9 中所见，由此可知 $f(F)$ 不含于 $L(p',q')$ 中。因此 $L(p',q')\subsetneq L(p,q)=L_0$，矛盾由此得出。

\par\medskip\noindent\textbf{步骤 11}\quad 现已知道 $r$ 是有理数。下面假定有理性定理中关于 $r$ 的分母的断言不成立。导出矛盾的手续几乎相同。

以 $\epsilon=1$ 应用步骤 3，即可如步骤 6 那样得到：存在任意大的 $(p,q)$，满足 $0<aq-rp<1$ 且
\[
P(p,q)=h^0(Y,pD_1+qD_2+\lceil A\rceil)\neq 0
\]
（否则 $P(x,y)$ 会过于频繁地消失）。于是与步骤 4 同样，$|pH+qa(K_X+\Delta)|\neq\varnothing$ 对所有 $(p,q)\in I$ 成立。
\end{proof}
% [Tutor Note: 全证的要点是把「$r$ 无理」或「分母过大」翻译成：带宽 $aq-rp$ 内落入了过多格点，从而 Hilbert 多项式的零点过于密集；引理 3.19 正是把这种密度转化为分母的上界。]
